\section*{摘要}

2022 年 FIA 引入"地面效应"空气动力学规则后，海豚跳（Porpoising）成为 F1 工程设计的关键挑战；
而空气动力学测试限制（ATR）对 CFD 仿真施加了严格的算力配额。
在"高成本仿真 $+$ 受限算力"的约束下，用代理模型替代昂贵仿真进行参数寻优成为可行路线。
鉴于所用数据集并非真实 CFD 输出，本文不以"优化真实赛车气动参数"为目标，
而研究一个更具普遍性的方法论问题：\textbf{在受限算力下，如何构建一个\emph{可信}
（Trustworthy）的代理模型优化框架——不仅求最优解，更要判断何时可以相信该最优解？}

方法上，本文构建"代理建模—不确定性估计—风险敏感优化—数据精炼—景观诊断"流程：
在 6 个模型谱系中选用 XGBoost 与 DeepEnsemble（R$^2$ 分别为 0.9999 与 0.9995）；
以 DeepEnsemble 的预测不确定性 $\sigma(x)$ 度量可信度，
并以 $\mu(x)-\lambda\sigma(x)$ 风险敏感目标抑制分布外（OOD）区域的过度外推；
通过不确定性引导的数据精炼检验代理的认知盲区。

研究发现：尽管代理预测精度极高，寻优却出现三项异常——预测越过物理上限、
不同优化器（PSO/GA/DE）的最优解高度接近、不同赛道场景的最优解高度趋同。
景观诊断表明，上述现象可由\textbf{优化景观退化}得到一种一致性的解释：
在目标变量饱和、关键变量高度共线的条件下，目标函数在高稳定区退化、区分度被抹平，
与优化结果趋同的现象一致。
据此，本文提出核心主张：\textbf{在代理模型驱动的优化中，数据覆盖度、
不确定性管理与优化景观结构，比优化器选择本身更决定最终优化质量。}
本文进而区分贡献（构建并验证了可信代理优化的诊断流程）
与发现（优化结果的趋同与景观退化、数据结构限制高度一致），
其方法论启示是：在代理驱动优化中，可信性的判断比单纯追求最优值更为关键。

\vspace{0.5cm}

\noindent
\textbf{关键词：}
可信优化；
代理模型；
不确定性感知；
粒子群优化；
景观诊断；
分布外外推；
计算智能；
F1 赛车气动稳定性
